\section{Scaling Laws y el futuro de los LLM}

\subsection{Scaling Laws: rendimientos decrecientes, pero lejos de terminar}

Ante el discurso reciente de que «Scaling Laws ya murió», Hassabis da una refutación clara: \textbf{los rendimientos de Scaling Laws están disminuyendo, pero están lejos de tocar un techo} («diminishing returns but not dead»).

\begin{importantbox}{El estado real de Scaling Laws}
\begin{itemize}
    \item Scaling Laws \textbf{no ha sido superada} (have NOT been hit)
    \item Los rendimientos están disminuyendo, pero la curva todavía baja
    \item La naturaleza de la competencia está cambiando: de «quién tiene más dinero» a «quién puede pensar algo nuevo»
\end{itemize}
\textit{«Rendimientos decrecientes, pero no muerta. El verdadero cambio es que la competencia pasa de 'quién tiene más dinero' a 'quién puede pensar algo nuevo'.» (Scaling: diminishing returns but not dead, competition shifting from ``who has more money'' to ``who can think of new things.'')}
\end{importantbox}

Este juicio implica que la época de simplemente acumular cómputo está quedando atrás, pero eso no significa que el escalamiento en sí deje de funcionar, sino que se entra en una nueva etapa que exige más pensamiento creativo. La innovación algorítmica, los nuevos paradigmas de entrenamiento y los diseños de arquitectura más eficientes se convertirán en el nuevo foco de la competencia.

\subsection{Cómputo: el mayor cuello de botella}

Hassabis señala con claridad que \textbf{el cómputo (Compute) es hoy el mayor cuello de botella} en el camino hacia la AGI. Aunque la competencia se está desplazando hacia la innovación algorítmica, la escala de los recursos de cómputo sigue siendo el factor clave que abre la brecha entre los laboratorios de IA de frontera.

\begin{warningbox}{Los LLM no se van a Commoditize}
Un malentendido común es que «al final los LLM se volverán un producto básico que cualquiera puede hacer». Hassabis se opone claramente a esto:
\begin{itemize}
    \item Los LLM no se van a comoditizar (will NOT be commoditized)
    \item La brecha entre las empresas de IA de frontera \textbf{no se está reduciendo, sino ampliando}
    \item Solo unas 4 empresas de frontera pueden seguir avanzando los modelos frontier
\end{itemize}
La razón es que el umbral de cómputo, datos, talento y know-how algorítmico que requieren los modelos de frontera está subiendo al mismo tiempo.
\end{warningbox}

\subsection{Los LLM son el componente base de la AGI}

Aunque los LLM tienen muchas limitaciones, Hassabis afirma con claridad que los LLM serán el \textbf{componente base} (foundational component) de la futura AGI, y no una tecnología de transición que será reemplazada. Lo importante no es abandonar los LLM, sino añadir sobre ellos las capacidades que les faltan: aprendizaje continuo, memoria, planeación y razonamiento consistente.

\begin{knowledgebox}{La ruta técnica de los LLM a la AGI}
Se puede entender el LLM como el «centro del lenguaje» de la AGI: aporta una capacidad potente de comprensión y generación del lenguaje, pero la AGI también necesita:
\begin{itemize}
    \item Un sistema de aprendizaje continuo (similar a la consolidación de memoria del hipocampo)
    \item Un módulo de planeación jerárquica (similar a la función ejecutiva de la corteza prefrontal)
    \item Un modelo del mundo (una simulación interna del mundo físico y social)
    \item Un motor de razonamiento consistente (que garantice la coherencia lógica entre dominios)
\end{itemize}
El LLM no será reemplazado, pero será potenciado y ampliado de forma considerable.
\end{knowledgebox}

\subsection{La brecha de percepción: sobrestimado a corto plazo, subestimado a largo plazo}

Hassabis señala un sesgo cognitivo importante:

\begin{importantbox}{La brecha de percepción sobre la IA}
\textit{«La IA está sobrevalorada a corto plazo y muy subestimada a largo plazo (a escala de 10 años).» (AI overhyped short-term, underappreciated long-term.)}

Esta brecha de percepción provoca dos problemas:
\begin{itemize}
    \item \textbf{Corto plazo}: las expectativas excesivas llevan a la decepción y al pánico de un «invierno de la IA»
    \item \textbf{Largo plazo}: subestimar la velocidad del cambio deja a la sociedad sin preparación suficiente
\end{itemize}
La actitud correcta es mantener el pragmatismo ante los avances recientes y el respeto ante el impacto de largo plazo.
\end{importantbox}

\subsection{Resumen de la sección}

Scaling Laws está lejos de terminar, pero el foco de la competencia está pasando de ser intensivo en capital a ser intensivo en innovación. Los LLM no se van a comoditizar, y la brecha entre los laboratorios de frontera se está ampliando, no reduciendo. El LLM seguirá siendo potenciado como componente base de la AGI. La percepción pública sobre la IA presenta un sesgo sistemático de «sobrestimación a corto plazo, subestimación a largo plazo».

%% ==================== Section 4 ====================
